\chapter{Introduzione}
\label{chap:intro}

Negli ultimi anni, i sistemi IoT si sono diffusi anche nell'ambito della protezione civile e della gestione delle emergenze, dove la disponibilità di informazioni tempestive e affidabili può incidere direttamente sull'efficacia dei soccorsi.
Il terremoto del 2016 che ha colpito l'Italia centrale, e in particolare l'area di Camerino, ha messo in evidenza questa esigenza. Infatti, nelle prime ore successive a un evento sismico, per i soccorritori è fondamentale sapere dove sono le persone, soprattutto quando potrebbero trovarsi in un luogo difficilmente accessibile o non essere in grado di muoversi. Da qui nasce la domanda alla base della tesi. Quale sensore, tra il PIR a infrarossi passivi e il radar a onde millimetriche (mmWave), rileva la presenza di una persona con maggiore affidabilità, e il radar può sostituire il PIR nei sistemi che oggi lo usano o deve affiancarlo?

Il capitolo è organizzato come segue. La Sezione~\ref{sec:dipme} descrive il progetto DIPME, contesto applicativo del lavoro, la Sezione~\ref{sec:obbiettivo} definisce gli obbiettivi della tesi e la Sezione~\ref{sec:struttura} ne illustra la struttura.


\section{Il progetto DIPME}
\label{sec:dipme}

Il confronto tra i due sensori trova uno scenario di applicazione concreto nel progetto DIPME, che ha come obbiettivo la mitigazione del rischio sismico attraverso l'integrazione di sensoristica IoT negli arredi degli ambienti scolastici \cite{slideDipme}. La tesi ne riprende il problema e ne valuta una possibile estensione. 
L'idea del progetto si articola su tre livelli:

\begin{itemize}
    \item \textbf{Arredi salva-vita.} Banchi e pareti sono progettati dalla Scuola di Architettura e Design di UNICAM per resistere al crollo del solaio e mantenere al di sotto di sé un volume in cui una persona possa trovare riparo (Figura~\ref{fig:arredo-render}). Non si tratta quindi di semplici mobili rinforzati, ma di elementi strutturali pensati per lavorare insieme e vincolati tra loro
    \cite{pietroni2021furniture}.

    \item \textbf{Nodi di sensing.} All'interno dell'arredo è integrata l'elettronica della piattaforma DIPME, sviluppata con partner industriali, fra cui AM Microsystems~\cite{slideDipme}, nell'ambito del progetto VITALITY~\cite{callisto2026dipme}. Ogni nodo, il DIPME-DEVICE, è un'unità autonoma alimentata a batteria e integra una radio LoRa, un sensore di presenza PIR e alcuni sensori ambientali per la misura di temperatura, umidità, pressione e anidride carbonica.

    \item \textbf{Raccolta e visualizzazione.} I nodi comunicano via LoRa con un coordinatore, che inoltra i dati a un gateway locale. LoRa permette di trasmettere a basse potenze anche su distanze dell'ordine di 3--5~km in ambiente urbano~\cite{slideDipme}. Il gateway, alimentato da un gruppo di continuità, coordina il comportamento dei nodi e trasmette i dati a una piattaforma cloud con cruscotti di monitoraggio, conservandoli quando la connessione si interrompe e reinviandoli al ripristino~\cite{callisto2026dipme}. La piattaforma SAFE, di cui DIPME estende l'approccio, prevede inoltre un sottosistema di localizzazione basato su droni e dispositivi mobili a supporto dei soccorritori~\cite{callisto2024safe,callisto2026dipme}.
\end{itemize}

\begin{figure}[htbp]
  \centering
  \includegraphics[width=0.9\textwidth]{Arredo_salva_vita}
  \caption{Il principio dell'arredo salva-vita. Il banco (a sinistra) e la parete attrezzata (a destra) conservano un volume di riparo sotto il carico del crollo. Il nodo di sensing, fissato sotto il piano, osserva l'occupante~\cite{slideDipme}.}
  \label{fig:arredo-render}
\end{figure}

Il sistema opera in due modalità, un tempo di \emph{pace}, dedicato al monitoraggio ordinario, e un tempo di \emph{guerra}, associato alla gestione dell'emergenza sismica. Il passaggio dalla prima alla seconda è comandato dal gateway al segnale di un accelerometro strutturale. Se il gateway non conferma i messaggi per tre ore, il nodo passa da solo al regime di emergenza. È in questa situazione che il rilevamento della presenza assume un'importanza particolare, perché da questo dato dipende la cadenza delle trasmissioni del nodo. Con una persona rilevata il nodo invia un messaggio al minuto, senza presenza uno ogni quindici minuti. Una persona sotto un banco che il sensore non rileva verrebbe quindi segnalata come assente, e il suo banco trasmetterebbe quindici volte meno spesso, proprio durante l'emergenza. DIPME è in esercizio in un dimostratore con 42 nodi installati nei banchi di due aule di una scuola secondaria delle Marche~\cite{callisto2026dipme}.

\subsection{Il limite del sensore adottato}
\label{sec:limite-pir}

Il PIR, la soluzione più diffusa per questa funzione grazie ai consumi e ai costi ridotti, ha però un limite che non dipende dalla qualità del componente, ma dal suo principio di funzionamento, discusso nel Capitolo~\ref{chap:sensori}. Il sensore risponde infatti alla variazione del flusso infrarosso, non al suo valore, e quindi non è in grado di rilevare una persona immobile. 

Nello scenario DIPME questo limite si presenta proprio nella situazione più critica, perché una persona rifugiata sotto un banco dopo un crollo è, tipicamente, ferma. Gli autori del progetto DIPME indicano, tra gli sviluppi futuri, l'estensione del sistema di sensing con radar mmWave a complemento del PIR, da valutare in termini di accuratezza, consumo, privacy e complessità \cite{callisto2026dipme}. Un radar a 24 GHz, sensibile ai micro-movimenti e in grado di misurare la distanza del bersaglio, può in linea di principio superare il limite del PIR.


\section{Obbiettivi}
\label{sec:obbiettivo}

La tesi risponde alla domanda posta in apertura con un confronto sperimentale fra un sensore PIR e due moduli radar mmWave a 24 GHz, l'\ldb\ e l'\ldc, basato su misure ripetibili che stabiliscono anche quanto sia ampio il divario fra le due tecnologie. Il risultato interessa direttamente il progetto DIPME.

Il lavoro è stato organizzato in sei punti:

\begin{itemize}
    \item \textbf{O1. Studiare i sensori}, per capirne il funzionamento, gli impieghi e i dati che producono.

    \item \textbf{O2. Analizzare i dati forniti}, individuando quali grandezze ciascun sensore mette a disposizione e con quale accuratezza.

    \item \textbf{O3. Confrontare sperimentalmente} PIR e mmWave con metriche quantitative.

    \item \textbf{O4. Valutare il consumo energetico}, a titolo informativo e sulla base dei datasheet, non essendo disponibile la strumentazione per misurarlo.

    \item \textbf{O5. Realizzare un'interfaccia web} che mostri i dati in tempo reale, li salvi e li esporti in formato CSV per l'analisi numerica.

    \item \textbf{O6. Definire un indice di vitalità}, calcolato a bordo del nodo, che oltre alla presenza misuri quanto una persona si muove.
\end{itemize}

Gli altri elementi del progetto DIPME (la radio LoRa, il gateway e la piattaforma cloud) e il sottosistema di localizzazione con i droni della piattaforma SAFE non sono oggetto della tesi e compaiono nella Sezione~\ref{sec:dipme} solo come contesto.

\section{Struttura della tesi}
\label{sec:struttura}

La tesi è organizzata come segue.

\begin{itemize}
    \item \textbf{I sensori di presenza.} Il Capitolo~\ref{chap:sensori} spiega come funzionano il PIR e il radar FMCW a 24 GHz e li colloca rispetto alle alternative (O1).
    
    \item \textbf{Analisi dei componenti hardware e dei dataset prodotti.} Il Capitolo~\ref{chap:moduli} presenta i tre sensori impiegati a partire dalla documentazione ufficiale, e la piattaforma di acquisizione su ESP32 con il formato dati comune a tutta la tesi (O2).
    
    \item \textbf{Metodologia di confronto.} Il Capitolo~\ref{chap:metodologia} fissa il protocollo, le metriche e le convenzioni con cui i sensori sono stati confrontati, e presenta l'interfaccia web servita direttamente dall'ESP32, senza alcuna rete esterna, realizzata a supporto della campagna (O5).

    \item \textbf{Confronto sperimentale PIR vs mmWave.} Il Capitolo~\ref{chap:confronto} è il cuore della tesi e raccoglie i risultati della campagna di misura, con ogni risultato accompagnato dalla sua dispersione su ripetizioni indipendenti (O3).

    \item \textbf{L'indice di vitalità.} Il Capitolo~\ref{chap:vitalita} definisce, tara e valida l'indice di vitalità, ne descrive il calcolo a bordo del microcontrollore e chiude con la stima del ritmo respiratorio, verificata contro un ritmo imposto (O6).

    \item \textbf{Consumo energetico.} Il Capitolo~\ref{chap:consumi} stima i consumi dei tre sensori e l'autonomia del nodo a batteria, e ne ricava i vincoli per il nodo del progetto (O4).

    \item \textbf{Conclusioni e sviluppi futuri.} Il Capitolo~\ref{chap:conclusioni} risponde alla domanda di partenza e indica come proseguire il lavoro.
\end{itemize}